\documentclass[12pt, a4paper]{article}

\usepackage[utf8]{inputenc}
\usepackage[margin=1in]{geometry}
\usepackage{graphicx}
\usepackage{hyperref}
\usepackage{listings}
\usepackage{xcolor}
\usepackage{booktabs}
\usepackage{amsmath}
\usepackage{enumitem}
\usepackage{float}

\definecolor{codegreen}{rgb}{0,0.6,0}
\definecolor{codegray}{rgb}{0.5,0.5,0.5}
\definecolor{codepurple}{rgb}{0.58,0,0.82}
\definecolor{backcolour}{rgb}{0.95,0.95,0.97}

\lstdefinestyle{pythonstyle}{
    backgroundcolor=\color{backcolour},
    commentstyle=\color{codegreen},
    keywordstyle=\color{blue},
    numberstyle=\tiny\color{codegray},
    stringstyle=\color{codepurple},
    basicstyle=\ttfamily\footnotesize,
    breaklines=true,
    frame=single,
    numbers=left,
    numbersep=5pt,
    showstringspaces=false,
    tabsize=4
}
\lstset{style=pythonstyle}

\title{\textbf{Vibe Forecaster} \\ \large Implementation Report — Test Implementation}
\author{ML Project — Exploratory Test}
\date{March 2026}

\begin{document}
\maketitle
\tableofcontents
\newpage

% ============================================================
\section{Introduction}

The \textbf{Vibe Forecaster} is a music mood prediction system that combines data from Last.fm (listening history) and Spotify (audio features) to:

\begin{enumerate}
    \item Cluster songs into meaningful ``vibe'' groups using K-Means clustering
    \item Train an LSTM neural network to predict the user's likely music mood at any given time
    \item Present results through an interactive Streamlit dashboard
\end{enumerate}

This implementation under \texttt{test1/} serves as an exploratory test to evaluate whether this approach yields meaningful and actionable predictions from real listening data.

% ============================================================
\section{System Architecture}

The system follows a sequential pipeline architecture:

\begin{verbatim}
Last.fm API  -->  Scrobble Data  -->  Merge + Feature  -->  K-Means
Spotify API  -->  Audio Features       Engineering         Clustering
                                                              |
                                                              v
Streamlit  <--  Prediction  <--  LSTM Training  <--  Clustered
Dashboard       Engine                                Dataset
\end{verbatim}

\subsection{Data Flow}
\begin{enumerate}
    \item \textbf{Extraction}: Scrobble history from Last.fm API (or existing CSV data)
    \item \textbf{Enrichment}: Spotify audio features for each unique track
    \item \textbf{Feature Engineering}: Temporal features + audio features merged
    \item \textbf{Clustering}: K-Means groups songs into vibe clusters
    \item \textbf{Training}: LSTM learns temporal patterns in vibe sequences
    \item \textbf{Prediction}: Real-time vibe forecasting
    \item \textbf{Visualization}: Interactive Streamlit dashboard
\end{enumerate}

% ============================================================
\section{Phase 0 — Environment Setup}

\subsection{Dependencies}
All required packages are listed in \texttt{requirements.txt}:

\begin{lstlisting}[language=bash]
pip install -r requirements.txt
\end{lstlisting}

Key libraries and their roles:
\begin{itemize}
    \item \texttt{pandas}, \texttt{numpy} — Data manipulation
    \item \texttt{spotipy} — Spotify Web API client
    \item \texttt{requests} — HTTP requests for Last.fm API
    \item \texttt{scikit-learn} — K-Means clustering, StandardScaler, PCA, metrics
    \item \texttt{tensorflow} — LSTM neural network
    \item \texttt{streamlit}, \texttt{plotly} — Interactive dashboard
    \item \texttt{matplotlib}, \texttt{seaborn} — Static visualizations
\end{itemize}

\subsection{API Credentials}
Copy \texttt{.env.template} to \texttt{.env} and populate with credentials from:
\begin{itemize}
    \item \href{https://developer.spotify.com/dashboard}{Spotify Developer Dashboard}
    \item \href{https://www.last.fm/api/account/create}{Last.fm API}
\end{itemize}

% ============================================================
\section{Phase 1 — Data Extraction}

\subsection{Script: \texttt{1\_extract\_lastfm.py}}
Fetches scrobble history from the Last.fm API using the \texttt{user.getrecenttracks} endpoint.

\textbf{Key design decisions:}
\begin{itemize}
    \item Fetches 200 tracks per page (API maximum) across 50 pages ($\approx$10,000 tracks)
    \item Adds 250ms delay between requests to respect rate limits
    \item Skips ``now playing'' tracks (no timestamp available)
    \item Converts Unix timestamps to datetime for downstream processing
\end{itemize}

\textbf{Output}: \texttt{data/scrobbles.csv} with columns: \texttt{artist}, \texttt{track}, \texttt{album}, \texttt{timestamp}, \texttt{datetime}

\subsection{Script: \texttt{1b\_convert\_existing\_data.py}}
Alternative path for users who already have Last.fm data exported as \texttt{fm\_data.csv}. Converts the existing format (Username, Artist, Track, Album, Date, Time) to the standard scrobbles format.

\subsection{Script: \texttt{2\_fetch\_audio\_features.py}}
Enriches each unique (artist, track) pair with Spotify's audio features via the Web API.

\textbf{Audio features retrieved:}
\begin{itemize}
    \item \textbf{valence} (0--1): Musical positiveness / happiness
    \item \textbf{energy} (0--1): Intensity and activity
    \item \textbf{danceability} (0--1): Suitability for dancing
    \item \textbf{acousticness} (0--1): Confidence the track is acoustic
    \item \textbf{instrumentalness} (0--1): Predicts whether a track has no vocals
    \item \textbf{tempo}: Estimated BPM
    \item \textbf{loudness}: Overall loudness in dB
\end{itemize}

\textbf{Rate limiting}: 100ms delay between requests to stay within Spotify's limits.

% ============================================================
\section{Phase 2 — Feature Engineering}

\subsection{Script: \texttt{3\_build\_dataset.py}}

Merges scrobble history with audio features and engineers temporal features:

\begin{table}[H]
\centering
\begin{tabular}{lll}
\toprule
\textbf{Feature} & \textbf{Type} & \textbf{Description} \\
\midrule
hour & Integer (0--23) & Hour of day when track was played \\
day\_of\_week & Integer (0--6) & 0=Monday, 6=Sunday \\
month & Integer (1--12) & Month of play \\
is\_weekend & Binary (0/1) & Whether Saturday or Sunday \\
valence & Float (0--1) & Musical positiveness \\
energy & Float (0--1) & Intensity \\
danceability & Float (0--1) & Dance suitability \\
acousticness & Float (0--1) & Acoustic confidence \\
instrumentalness & Float (0--1) & Instrumental prediction \\
tempo & Float & BPM \\
loudness & Float (dB) & Overall loudness \\
\bottomrule
\end{tabular}
\caption{Master dataset feature descriptions}
\end{table}

\textbf{Output}: \texttt{data/master\_dataset.csv}

% ============================================================
\section{Phase 3 — Exploratory Data Analysis}

\subsection{Script: \texttt{4\_eda.py}}

Generates five visualizations to understand listening patterns before modeling:

\begin{enumerate}
    \item \textbf{Correlation Heatmap}: Shows relationships between audio features (e.g., energy vs. acousticness are typically negatively correlated)
    \item \textbf{Hourly Trends}: Average energy and valence by hour of day — reveals daily mood cycles
    \item \textbf{Energy Heatmap}: Day-of-week $\times$ hour matrix showing when high-energy listening peaks occur
    \item \textbf{Feature Distributions}: Histograms of each audio feature to check for skew or bimodality
    \item \textbf{Top Artists}: Most-listened artists by scrobble count
\end{enumerate}

These plots inform clustering decisions (e.g., whether tempo has enough variance to be useful as a clustering feature).

% ============================================================
\section{Phase 4 — Vibe Clustering (K-Means)}

\subsection{Script: \texttt{5\_clustering.py}}

\subsubsection{Feature Scaling}
All 7 audio features are standardized using \texttt{StandardScaler} (zero mean, unit variance) before clustering, since K-Means is distance-based and sensitive to feature scales.

\subsubsection{Choosing $k$ — Elbow Method}
We evaluate $k \in \{2, 3, \ldots, 9\}$ and plot inertia (within-cluster sum of squares) vs. $k$. The ``elbow'' point where diminishing returns begin suggests the optimal number of clusters. We default to $k=5$, which typically produces interpretable vibes.

\subsubsection{Cluster Naming}
Clusters are automatically named based on their centroid feature profiles using heuristic rules:

\begin{itemize}
    \item \textbf{Party Mode}: High energy + high valence + high danceability
    \item \textbf{Peak Energy}: High energy + high valence (not dance-focused)
    \item \textbf{Chill Vibes}: High acousticness + low energy
    \item \textbf{Deep Focus}: High instrumentalness
    \item \textbf{Wind-down}: Low valence + low energy
\end{itemize}

\subsubsection{Visualization}
PCA reduces the 7-dimensional feature space to 2 dimensions for scatter plot visualization, with each cluster colored differently. The explained variance ratio is displayed on axis labels to show how much information the projection retains.

\textbf{Outputs}:
\begin{itemize}
    \item \texttt{data/clustered\_dataset.csv} — Master dataset with vibe labels
    \item \texttt{models/kmeans.pkl} — Trained K-Means model
    \item \texttt{models/scaler.pkl} — Fitted StandardScaler
\end{itemize}

% ============================================================
\section{Phase 5 — LSTM Vibe Forecasting}

\subsection{Script: \texttt{6\_train\_lstm.py}}

\subsubsection{Sequence Construction}
The LSTM operates on fixed-length sequences of historical listening events. For each position $i$ in the chronologically-sorted dataset, we construct:

$$X_i = \begin{bmatrix} \text{vibe\_id}_{i} & \text{hour}_{i} & \text{dow}_{i} & \text{weekend}_{i} \\ \vdots & \vdots & \vdots & \vdots \\ \text{vibe\_id}_{i+L-1} & \text{hour}_{i+L-1} & \text{dow}_{i+L-1} & \text{weekend}_{i+L-1} \end{bmatrix}$$

$$y_i = \text{vibe\_id}_{i+L}$$

where $L = 20$ is the sequence length (configurable via \texttt{SEQ\_LEN}).

\subsubsection{Model Architecture}

\begin{verbatim}
Input (SEQ_LEN × 4)
    |
LSTM(64, return_sequences=True)
    |
Dropout(0.2)
    |
LSTM(32)
    |
Dropout(0.2)
    |
Dense(NUM_CLASSES, softmax)
\end{verbatim}

\begin{itemize}
    \item \textbf{Two-layer LSTM}: First layer returns sequences to allow the second LSTM layer to process temporal dependencies at a higher abstraction level
    \item \textbf{Dropout (20\%)}: Regularization to prevent overfitting on the training sequences
    \item \textbf{Softmax output}: Produces probability distribution over all vibe classes
    \item \textbf{Loss}: Sparse categorical cross-entropy (labels are integers, not one-hot)
    \item \textbf{Optimizer}: Adam with default learning rate
\end{itemize}

\subsubsection{Training Strategy}
\begin{itemize}
    \item \textbf{No shuffle}: Data is split 80/20 chronologically (not randomly) to respect temporal ordering — the model should predict future vibes from past data, not interpolate
    \item \textbf{30 epochs, batch size 32}: Standard starting configuration
    \item Training and validation curves are plotted to monitor overfitting
\end{itemize}

\textbf{Outputs}:
\begin{itemize}
    \item \texttt{models/lstm\_vibe.h5} — Trained LSTM model
    \item \texttt{models/label\_encoder.pkl} — Vibe label encoder
\end{itemize}

% ============================================================
\section{Phase 6 — Prediction Engine}

\subsection{Script: \texttt{7\_predict.py}}

The prediction pipeline at inference time:

\begin{enumerate}
    \item Load the last $L=20$ scrobbles from the clustered dataset
    \item Encode their vibe labels to integer IDs
    \item Pack each into a feature vector $[\text{vibe\_id}, \text{hour}, \text{dow}, \text{weekend}]$
    \item Feed the sequence into the LSTM
    \item Output: probability distribution over all vibe classes
    \item Report the top vibe and its confidence percentage
\end{enumerate}

The output includes a visual bar chart in the terminal showing all vibe probabilities, making it easy to see the full distribution.

% ============================================================
\section{Phase 7 — Streamlit Dashboard}

\subsection{Script: \texttt{app.py}}

The interactive dashboard provides four main views:

\begin{enumerate}
    \item \textbf{Current Vibe Prediction}: Real-time prediction with confidence metric and probability bar chart
    \item \textbf{Weekly Heatmap}: Day $\times$ hour matrix showing the most common vibe cluster at each time slot (Plotly \texttt{imshow})
    \item \textbf{Cluster Profiles}: Grouped bar chart comparing mean audio features across all vibe clusters
    \item \textbf{Distribution Views}: Pie chart of overall vibe distribution and hourly listening activity
\end{enumerate}

Models are cached using \texttt{@st.cache\_resource} to avoid reloading on each interaction.

\textbf{Run}: \texttt{streamlit run app.py}

% ============================================================
\section{Project Structure}

\begin{verbatim}
test1/
├── .env.template          # API credential template
├── requirements.txt       # Python dependencies
├── fm_data.csv           # Existing Last.fm export (input)
├── 1_extract_lastfm.py   # Phase 1: Fetch from Last.fm API
├── 1b_convert_existing_data.py  # Phase 1: Convert existing CSV
├── 2_fetch_audio_features.py    # Phase 1: Spotify enrichment
├── 3_build_dataset.py    # Phase 2: Feature engineering
├── 4_eda.py              # Phase 3: Exploratory analysis
├── 5_clustering.py       # Phase 4: K-Means clustering
├── 6_train_lstm.py       # Phase 5: LSTM training
├── 7_predict.py          # Phase 6: CLI prediction
├── app.py                # Phase 7: Streamlit dashboard
├── data/                 # Generated data files
│   ├── scrobbles.csv
│   ├── audio_features.csv
│   ├── master_dataset.csv
│   └── clustered_dataset.csv
├── models/               # Trained model artifacts
│   ├── kmeans.pkl
│   ├── scaler.pkl
│   ├── label_encoder.pkl
│   └── lstm_vibe.h5
├── plots/                # Generated visualizations
│   ├── correlation.png
│   ├── hourly_trends.png
│   ├── energy_heatmap.png
│   ├── feature_distributions.png
│   ├── top_artists.png
│   ├── elbow.png
│   ├── clusters_pca.png
│   ├── cluster_profiles.png
│   └── training_history.png
└── docs/
    └── implementation_report.tex
\end{verbatim}

% ============================================================
\section{Execution Order}

Scripts must be run in sequence, as each depends on the output of previous steps:

\begin{enumerate}
    \item \texttt{1\_extract\_lastfm.py} (or \texttt{1b\_convert\_existing\_data.py})
    \item \texttt{2\_fetch\_audio\_features.py}
    \item \texttt{3\_build\_dataset.py}
    \item \texttt{4\_eda.py}
    \item \texttt{5\_clustering.py}
    \item \texttt{6\_train\_lstm.py}
    \item \texttt{7\_predict.py} (standalone) or \texttt{streamlit run app.py}
\end{enumerate}

% ============================================================
\section{Known Limitations \& Future Work}

\subsection{Current Limitations}
\begin{itemize}
    \item \textbf{Spotify API dependency}: Audio features require Spotify to find matching tracks — unmatched tracks are dropped, reducing dataset size
    \item \textbf{Static prediction context}: The model uses the last 20 historical scrobbles, not real-time streaming data
    \item \textbf{Cluster naming heuristic}: Automatic naming may not always produce intuitive labels — manual inspection of centroids recommended
    \item \textbf{Data sparsity}: LSTM performance degrades with fewer than $\sim$5,000 scrobbles (sequence diversity is insufficient)
\end{itemize}

\subsection{Potential Improvements}
\begin{itemize}
    \item Add weather/calendar context features for richer temporal modeling
    \item Use Transformer architecture instead of LSTM for longer-range dependencies
    \item Implement online learning to update the model as new scrobbles arrive
    \item Add song recommendation within the predicted vibe cluster
    \item Deploy as a persistent web service with scheduled data updates
\end{itemize}

% ============================================================
\section{Conclusion}

This test implementation demonstrates a complete end-to-end pipeline from raw listening data to interactive vibe prediction. The K-Means clustering successfully separates songs into interpretable mood categories based on Spotify's audio features, and the LSTM model captures temporal patterns in listening behavior.

The Streamlit dashboard provides an accessible interface for exploring predictions and listening patterns. This exploratory test validates the overall approach and identifies areas for refinement before a production deployment.

\end{document}
